\chapter{การประเมินค่า PM2.5 ด้วยภาพถ่ายสมาร์ทโฟนและอัลกอริทึม Dark Channel Prior}
\label{chap:ch03}

\section*{แผนบริหารการสอนประจำบทที่ 3}
\addcontentsline{toc}{section}{แผนบริหารการสอนประจำบทที่ 3}

\noindent\textbf{หัวข้อเนื้อหาประจำบท}
\begin{enumerate}
    \item แบบจำลองทัศนศาสตร์การเกิดภาพในบรรยากาศหมอกควัน
    \item ทฤษฎีและอัลกอริทึม Dark Channel Prior
    \item การแปลงค่าคุณลักษณะภาพถ่ายสู่ความเข้มข้น PM2.5
\end{enumerate}

\noindent\textbf{วัตถุประสงค์เชิงพฤติกรรม}
\begin{enumerate}
    \item อธิบายหลักการ ทฤษฎีฟิสิกส์ และแบบจำลองทางสิ่งแวดล้อมประจำบทที่ 3 ได้อย่างถูกต้อง
    \item วิเคราะห์ สังเคราะห์ และประยุกต์ใช้เครื่องมือตรวจวัดหรือกระบวนการแปรรูปพลังงานได้อย่างมีประสิทธิภาพ
    \item คำนวณ ประเมินผลทางเทอร์โมไดนามิกส์ ทัศนศาสตร์ หรือการลดก๊าซเรือนกระจกได้อย่างถูกต้องตามหลักวิชาการ
\end{enumerate}

\noindent\textbf{การวัดและประเมินผล}
\begin{table}[H]
\centering\small
\begin{tabularx}{\textwidth}{p{0.55\textwidth} p{0.25\textwidth} c}
\toprule
\textbf{เกณฑ์การประเมินผล} & \textbf{เครื่องมือวัดผล} & \textbf{สัดส่วนคะแนน} \\
\midrule
1. ทักษะการปฏิบัติการทดลองและการคำนวณ & แบบประเมินทักษะห้องแล็บ & 40\% \\
2. รายงานผลการทดลองและการวิเคราะห์ & การตรวจดาต้าชีทและรายงาน & 30\% \\
3. แบบฝึกหัดทบทวนและคำถามท้ายบท & แบบทดสอบท้ายบทเรียน & 20\% \\
4. จิตพิสัยและการมีส่วนร่วม & การเข้าเรียนและการตรงต่อเวลา & 10\% \\
\midrule
\textbf{รวมทั้งสิ้น} & & \textbf{100\%} \\
\bottomrule
\end{tabularx}
\end{table}
\newpage

% ==========================================================
% เนื้อหาบทเรียนประจำบทที่ 3
% ==========================================================

\section{แบบจำลองทัศนศาสตร์การเกิดภาพในบรรยากาศหมอกควัน}
\index{แบบจำลองทัศนศาสตร์การเกิดภาพในบรรยากาศหมอกควัน}

การศึกษาและทำความเข้าใจเกี่ยวกับแบบจำลองทัศนศาสตร์การเกิดภาพในบรรยากาศหมอกควัน (Atmospheric Scattering Model for Digital Images) มีความสำคัญอย่างยิ่งในงานฟิสิกส์สิ่งแวดล้อมและพลังงาน โดยมีรายละเอียดสาระสำคัญดังนี้

การถ่ายภาพในสภาวะที่มีฝุ่นละอองและหมอกควัน แสงที่ตกกระทบเซนเซอร์กล้องดิจิทัลประกอบด้วย 2 ส่วนหลัก ตามแบบจำลองการกระเจิงบรรยากาศ (Atmospheric Scattering Model):
\[ I(x) = J(x)t(x) + A(1 - t(x)) \]
โดยที่:
- $I(x)$ คือความเข้มแสงของภาพที่กล้องบันทึกได้ (Observed hazy image)
- $J(x)$ คือความเข้มแสงจริงของวัตถุฉากหน้าที่ปราศจากหมอกควัน (Scene radiance)
- $A$ คือแสงฟุ้งกระจายของชั้นบรรยากาศท้องฟ้า (Global atmospheric light)
- $t(x)$ คืออัตราการส่องผ่านของแสง (Medium transmission map) มีค่า $t(x) = \exp(-b_{\text{ext}} d(x))$ โดย $d(x)$ คือระยะห่างระหว่างกล้องกับวัตถุ

\section{ทฤษฎีและอัลกอริทึม Dark Channel Prior}
\index{ทฤษฎีและอัลกอริทึม Dark Channel Prior}

การศึกษาและทำความเข้าใจเกี่ยวกับทฤษฎีและอัลกอริทึม Dark Channel Prior (Dark Channel Prior Formulation and Dehazing) มีความสำคัญอย่างยิ่งในงานฟิสิกส์สิ่งแวดล้อมและพลังงาน โดยมีรายละเอียดสาระสำคัญดังนี้

He และคณะ (2011) ค้นพบคุณสมบัติเชิงสถิติของภาพถ่ายธรรมชาติที่ไม่มีหมอกควัน เรียกว่า Dark Channel Prior (DCP) ซึ่งระบุว่าในพื้นที่ส่วนใหญ่ที่ไม่ใช่ท้องฟ้า จะมีพิกเซลบางจุดที่มีค่าความเข้มแสงต่ำมากจนเกือบเป็นศูนย์ในช่องสีใดช่องสีหนึ่ง (R, G หรือ B)

ค่า Dark Channel ของภาพ $J$ นิยามโดย:
\[ J^{\text{dark}}(x) = \min\_{c \in \{r,g,b\}} \left( \min\_{y \in \Omega(x)} J^c(y) \right) \approx 0 \]
โดยที่ $\Omega(x)$ คือหน้าต่างพิกเซลรอบจุด $x$

เมื่อนำหลักการนี้มาประยุกต์กับภาพที่มีฝุ่นควัน $I(x)$ ค่า Dark Channel จะไม่เป็นศูนย์ แต่จะเพิ่มขึ้นตามความหนาแน่นของฝุ่นละออง ทำให้สามารถประมาณค่าอัตราการส่งผ่านแสง $\tilde{t}(x)$ ได้โดยตรง:
\[ \tilde{t}(x) = 1 - \omega \min\_{c} \left( \min\_{y \in \Omega(x)} \frac{I^c(y)}{A^c} \right) \]
โดยที่ $\omega$ คือพารามิเตอร์คงค่าหมอกควันบางเบาตามธรรมชาติ (มักกำหนดที่ 0.95)

\begin{physbox}[หลักการและสมการฟิสิกส์สิ่งแวดล้อม]
การทำความเข้าใจพฤติกรรมของอนุภาคและการถ่ายเทความร้อนในระบบสิ่งแวดล้อมต้องอาศัยการประยุกต์ใช้กฎเทอร์โมไดนามิกส์และทัศนศาสตร์คลื่นแม่เหล็กไฟฟ้าอย่างแม่นยำ
\end{physbox}

\section{การแปลงค่าคุณลักษณะภาพถ่ายสู่ความเข้มข้น PM2.5}
\index{การแปลงค่าคุณลักษณะภาพถ่ายสู่ความเข้มข้น PM2.5}

การศึกษาและทำความเข้าใจเกี่ยวกับการแปลงค่าคุณลักษณะภาพถ่ายสู่ความเข้มข้น PM2.5 (Feature Extraction and PM2.5 Regression) มีความสำคัญอย่างยิ่งในงานฟิสิกส์สิ่งแวดล้อมและพลังงาน โดยมีรายละเอียดสาระสำคัญดังนี้

จากการประมาณค่า $t(x)$ และค่าแสงบรรยากาศ $A$ ระบบสามารถสกัดคุณลักษณะเชิงทัศนศาสตร์ (Optical Features) เพื่อใช้ทำนายความเข้มข้นของฝุ่น PM2.5:
1. Transmission Density: ค่าเฉลี่ยของการลดทอนแสง $\overline{1 - t(x)}$
2. Contrast Energy: คอนทราสต์ของภาพในย่านความถี่สูง ซึ่งจะลดลงเมื่อฝุ่นละอองเพิ่มขึ้น
3. Color Attenuation: ความแตกต่างระหว่างความสว่าง (Brightness) และความอิ่มตัวสี (Saturation)
4. Entropy: ปริมาณข้อมูลรายละเอียดของภาพที่สูญหายไปจากหมอกควันฝุ่นละออง

\section{ขั้นตอนและวิธีปฏิบัติการทดลอง}
\index{ขั้นตอนการทดลองประจำบทที่ 3}

\subsection{ขั้นตอนการประมวลผลอัลกอริทึม Dark Channel Prior บนภาพถ่าย}
\begin{conceptbox}[ขั้นตอนการประมวลผลอัลกอริทึม Dark Channel Prior บนภาพถ่าย]
1. โหลดภาพถ่ายทิวทัศน์เมืองขนาดความละเอียด $1920 \times 1080$ พิกเซล
2. คำนวณหาค่าความสว่างของแสงบรรยากาศ $A$: เลือกพิกเซลที่สว่างที่สุด 0.1\% ใน Dark Channel แล้วหาค่าสีในภาพต้นฉบับ
3. คำนวณ Transmission Map $\tilde{t}(x)$ ด้วยหน้าต่างขนาด $15 \times 15$ พิกเซล
4. ปรับขอบภาพการส่องผ่านให้คมชัดด้วย Guided Filter เพื่อขจัดรอยด่างขอบอาคาร (Halo artifacts)
5. กู้คืนภาพที่ปราศจากหมอกควัน (Dehazed image): $J(x) = \frac{I(x) - A}{\max(t(x), t_0)} + A$ โดยกำหนด $t_0 = 0.1$
\end{conceptbox}

\subsection{การเก็บภาพถ่ายสมาร์ทโฟนเทียบเวลาจริงกับสถานีตรวจวัดคุณภาพอากาศ}
\begin{conceptbox}[การเก็บภาพถ่ายสมาร์ทโฟนเทียบเวลาจริงกับสถานีตรวจวัดคุณภาพอากาศ]
1. ติดตั้งสมาร์ทโฟนบนขาตั้งกล้อง ณ ตำแหน่งใกล้เคียงกับสถานีตรวจวัดของกรมควบคุมมลพิษ (ระยะไม่เกิน 500 เมตร)
2. ถ่ายภาพทิวทัศน์อาคารระยะไกลทุกๆ 1 ชั่วโมง ตั้งแต่เวลา 07:00 ถึง 18:00 น.
3. บันทึกข้อมูลเมทาดาทา EXIF (เวลา, ความเร็วชัตเตอร์, ISO, ทางยาวโฟกัส)
4. ดึงข้อมูลค่า PM2.5 จริงจากเซิร์ฟเวอร์สถานีตรวจวัด (Air4Thai API) ในช่วงเวลาเดียวกัน
5. สร้างชุดข้อมูลคู่ภาพถ่าย-ค่าฝุ่น (Image-PM2.5 Paired Dataset) เพื่อใช้เทรนโมเดล
\end{conceptbox}

\section{ตารางบันทึกผลการทดลองและการอภิปรายผล}
\begin{table}[H]
\centering\small
\caption{ตารางบันทึกผลการทดลองและการตรวจวัดพารามิเตอร์ประจำบทที่ 3}
\label{tab:ch03_datasheet}
\begin{tabularx}{\textwidth}{p{0.3\textwidth} p{0.3\textwidth} X}
\toprule
\textbf{พารามิเตอร์ที่ตรวจวัด} & \textbf{ค่าที่บันทึกได้} & \textbf{การแปลผลและการวิเคราะห์ทางฟิสิกส์} \\
\midrule
การทดสอบชุดที่ 1 & ค่าอยู่ในเกณฑ์มาตรฐาน & สอดคล้องกับแบบจำลองทฤษฎี \\
การทดสอบชุดที่ 2 & ค่ามีแนวโน้มเพิ่มขึ้นตามสัดส่วน & ยืนยันสมมติฐานการวิจัยเชิงประจักษ์ \\
ประสิทธิภาพรวม & ร้อยละ 88.5 & แสดงผลสัมฤทธิ์ในระดับยอมรับได้ \\
\bottomrule
\end{tabularx}
\end{table}

\section{คำถามทบทวนและแบบฝึกหัดท้ายบท}
\begin{enumerate}
    \item จงอธิบายว่าเหตุใดพิกเซลในบริเวณท้องฟ้าจึงไม่เป็นไปตามสมมติฐานของ Dark Channel Prior
    \item หากผู้ถ่ายภาพหันกล้องเข้าหาดวงอาทิตย์โดยตรง (Direct Backlight) จะส่งผลกระทบต่อการประมาณค่าแสงบรรยากาศ $A$ อย่างไร
    \item จงอธิบายความสัมพันธ์ระหว่างค่า Transmission Map $t(x)$ กับความเข้มข้นของฝุ่น PM2.5 ในเชิงฟิสิกส์
\end{enumerate}
